\section{Introduction}
\label{sec:introduction}

A trusted research environment (TRE) lets researchers analyze sensitive data while controlling what they can take out~\cite{fivesafes,hdruk2020tre}. Statistical disclosure control (SDC) suppresses small cells, rounds counts and withholds tables so that no individual can be identified from a released result~\cite{hundepool2012,griffiths2019handbook}. These rules protect privacy, but they can also remove the evidence that an analysis is wrong. A missing count of participants at each stage can hide that some were dropped. Withheld subgroup tables can hide that a pooled comparison reverses within every subgroup. A released output can therefore satisfy every disclosure rule and still support a wrong conclusion.

The question becomes practical when a language model does the analysis. A companion system places a model analyst behind a deterministic gateway that restricts its queries and audits every release~\cite{wade2026safetre}. That boundary governs what the model can disclose. It says nothing about whether the model reads the released aggregates correctly. Before a TRE admits such an analyst, it needs evidence on two things: how often the model draws a wrong conclusion, and how well a reviewer can tell when it has.

Existing benchmarks score data-analysis agents against reference answers~\cite{gu2024blade,majumder2024discoverybench,zhu2024statqa,liu2024qrdata,egg2025dabstep,jing2024dsbench,mitchener2025bixbench}. None measures how disclosure rules affect a reviewer who sees only the released output, and an analyst's benchmark score does not show that a second model can find its mistakes. These are separate measurements, with different denominators and different ways to fail.

OPACITY has two studies. The first asks whether a model can review an analysis under increasingly strict disclosure rules: which analytic errors remain detectable as more is withheld, and how much detection rests on the wording of the interpretation rather than on the numbers. The second evaluates a model as an analyst before admission: how the prompt, the reasoning budget and an independent checker change the rate of wrong conclusions released, and whether a configuration chosen on development tasks holds up on tasks it has never seen. A companion paper uses the same tasks and the configuration recommended here to test for bias and deception~\cite{wade2026vetting}.

The paper makes three contributions. It introduces a matched benchmark of eight analytic error classes under graded disclosure control, with controls for the wording of the interpretation. It evaluates model analysts, alone and followed by a checker, under selection rules fixed before two sealed task sets were opened. It measures checkers on analysts' real write-ups as well as on planted errors, and shows that detection of planted errors overstates detection of real ones. The accompanying materials preserve the prompts, model revisions, verdicts and selection procedures needed to repeat the evaluations on another model. Every evaluation reported here is by machine; the pre-registered human reviewer arm has not yet run.
